\textbf{Phần 1. Trắc nghiệm lựa chọn}
\setcounterpageref{ex}{0}

\begin{ex} %[4X9N51]
Tung một đồng xu 10 lần, có 6 lần xuất hiện mặt ngửa. Số lần xuất hiện mặt sấp là
\choice
{$3$ lần}
{$5$ lần}
{$6$ lần}
{\True $4$ lần}
\loigiai{
$10 - 6 = 4$ lần.\\
\textbf{Đáp án của câu hỏi là D}
}
\end{ex}

\begin{ex} %[4X9N51]
Khi gieo một con xúc xắc 6 mặt, các sự kiện có thể xảy ra ở mặt trên là xuất hiện số chấm từ
\choice
{$1$ đến $5$}
{\True $1$ đến $6$}
{$0$ đến $6$}
{$1$ đến $10$}
\loigiai{
Con xúc xắc có 6 mặt ghi từ 1 đến 6 chấm.\\
\textbf{Đáp án của câu hỏi là B}
}
\end{ex}

\begin{ex} %[4X9H51]
Gieo một con xúc xắc 20 lần, thấy có 4 lần xuất hiện mặt 6 chấm. Số lần xuất hiện các mặt không phải 6 chấm là
\choice
{$14$ lần}
{$15$ lần}
{\True $16$ lần}
{$18$ lần}
\loigiai{
$20 - 4 = 16$ lần.\\
\textbf{Đáp án của câu hỏi là C}
}
\end{ex}

\begin{ex} %[4X9H51]
Trong một hộp có 3 quả bóng đỏ và 2 quả bóng xanh. Lấy ngẫu nhiên 1 quả bóng. Sự kiện nào sau đây KHÔNG THỂ xảy ra?
\choice
{\True Lấy được bóng vàng}
{Lấy được bóng đỏ}
{Lấy được bóng xanh}
{Lấy được 1 quả bóng}
\loigiai{
Trong hộp không có bóng vàng nên không thể lấy được bóng vàng.\\
\textbf{Đáp án của câu hỏi là A}
}
\end{ex}

\begin{ex} %[4X9H51]
Tung đồng xu 15 lần, ghi lại kết quả: N, S, N, N, S, S, N, S, N, N, S, N, S, N, N (N: Ngửa, S: Sấp). Số lần xuất hiện mặt ngửa là
\choice
{$8$ lần}
{$10$ lần}
{$6$ lần}
{\True $9$ lần}
\loigiai{
Đếm số chữ N: có 9 lần.\\
\textbf{Đáp án của câu hỏi là D}
}
\end{ex}

\begin{ex} %[4X9H51]
Từ kết quả ở câu trên, số lần xuất hiện mặt sấp là
\choice
{$5$ lần}
{\True $6$ lần}
{$7$ lần}
{$8$ lần}
\loigiai{
$15 - 9 = 6$ lần.\\
\textbf{Đáp án của câu hỏi là B}
}
\end{ex}

\begin{ex} %[4X9V51]
Bạn Nam quay một chiếc vòng quay có 4 ô màu: Xanh, Đỏ, Vàng, Trắng 30 lần. Kim dừng lại ở ô màu Đỏ 8 lần, ô màu Xanh 10 lần, ô màu Vàng 7 lần. Số lần kim dừng lại ở ô màu Trắng là
\choice
{\True $5$ lần}
{$4$ lần}
{$6$ lần}
{$7$ lần}
\loigiai{
$30 - (8 + 10 + 7) = 30 - 25 = 5$ lần.\\
\textbf{Đáp án của câu hỏi là A}
}
\end{ex}

\begin{ex} %[4X9V51]
Trong trò chơi bốc thăm trúng thưởng, có 50 lá thăm gồm 10 lá thăm may mắn trúng thưởng và 40 lá thăm chúc bạn may mắn lần sau. Bạn Lan bốc ngẫu nhiên 1 lá thăm. Sự kiện nào có khả năng xảy ra cao hơn?
\choice
{Bốc được lá thăm trúng thưởng}
{\True Bốc được lá thăm chúc bạn may mắn lần sau}
{Hai sự kiện có khả năng như nhau}
{Chắc chắn trúng thưởng}
\loigiai{
40 lá thăm không trúng $>$ 10 lá thăm trúng nên khả năng bốc thăm không trúng cao hơn.\\
\textbf{Đáp án của câu hỏi là B}
}
\end{ex}

\begin{ex} %[4X9V51]
Gieo một con xúc xắc 50 lần, ghi nhận số lần xuất hiện mặt chẵn (2, 4, 6) là 28 lần. Số lần xuất hiện mặt lẻ (1, 3, 5) là
\choice
{\True $22$ lần}
{$20$ lần}
{$24$ lần}
{$26$ lần}
\loigiai{
$50 - 28 = 22$ lần.\\
\textbf{Đáp án của câu hỏi là A}
}
\end{ex}

\begin{ex} %[4X9V51]
Một túi kẹo có 10 cái kẹo dâu và 10 cái kẹo cam. Bạn Bình lấy ngẫu nhiên 1 cái kẹo. Khả năng lấy được kẹo dâu và kẹo cam là
\choice
{Kẹo dâu cao hơn}
{Kẹo cam cao hơn}
{\True Như nhau (bằng nhau)}
{Không thể lấy được kẹo dâu}
\loigiai{
Số lượng kẹo mỗi loại bằng nhau (10 cái) nên khả năng là như nhau.\\
\textbf{Đáp án của câu hỏi là C}
}
\end{ex}

\begin{ex} %[4X9V51]
Khi thực hiện 40 lần tung đồng xu, số lần xuất hiện mặt ngửa gấp 3 lần số lần xuất hiện mặt sấp. Số lần xuất hiện mặt ngửa là
\choice
{$20$ lần}
{$25$ lần}
{$35$ lần}
{\True $30$ lần}
\loigiai{
Tổng số phần: $1 + 3 = 4$ phần.\\
Mặt ngửa: $40 : 4 \times 3 = 30$ lần.\\
\textbf{Đáp án của câu hỏi là D}
}
\end{ex}

\begin{ex} %[4X9C51]
Trong hộp kín có các viên bi xanh và đỏ. Bạn An lấy ngẫu nhiên 1 viên bi, ghi lại màu rồi trả lại hộp, lặp lại 100 lần. Kết quả: 70 lần lấy được bi đỏ, 30 lần lấy được bi xanh. Dự đoán nào sau đây là hợp lí nhất?
\choice
{Số bi xanh trong hộp nhiều hơn số bi đỏ}
{Trong hộp chỉ có toàn bi đỏ}
{\True Số bi đỏ trong hộp nhiều hơn số bi xanh}
{Trong hộp có số bi xanh bằng bi đỏ}
\loigiai{
Tần số xuất hiện bi đỏ cao hơn nhiều (70 lần so với 30 lần) nên số bi đỏ trong hộp nhiều hơn.\\
\textbf{Đáp án của câu hỏi là C}
}
\end{ex}

\textbf{Phần 2. Trắc nghiệm đúng sai}
\setcounterpageref{ex}{0}

\begin{ex} %[4X9H51]
Gieo một con xúc xắc 30 lần: mặt 1 chấm xuất hiện 5 lần, mặt 2 chấm xuất hiện 6 lần, mặt 3 chấm xuất hiện 4 lần, mặt 4 chấm xuất hiện 5 lần, mặt 5 chấm xuất hiện 7 lần, mặt 6 chấm xuất hiện 3 lần. Xét tính đúng sai:
\choiceTF
{\True Mặt 5 chấm xuất hiện nhiều lần nhất ($7$ lần).}
{\True Mặt 6 chấm xuất hiện ít lần nhất ($3$ lần).}
{Tổng số lần xuất hiện mặt có số chấm là số chẵn là $15$ lần.}
{\True Tổng số lần xuất hiện mặt có số chấm là số chẵn là $14$ lần ($6 + 5 + 3 = 14$).}
\loigiai{
\begin{itemize}
    \item Mặt 5 chấm: 7 lần (nhiều nhất). Đúng.
    \item Mặt 6 chấm: 3 lần (ít nhất). Đúng.
    \item Số chẵn (2, 4, 6): $6 + 5 + 3 = 14$ lần $\neq 15$ lần. Ý 3 Sai, Ý 4 Đúng.
\end{itemize}
\begin{itemize}
    \item a) \textbf{Đúng}.
    \item b) \textbf{Đúng}.
    \item c) \textbf{Sai}.
    \item d) \textbf{Đúng}.
\end{itemize}
}
\end{ex}

\begin{ex} %[4X9V51]
Trong hộp có 5 viên bi xanh, 3 viên bi đỏ và 2 viên bi vàng. Lấy ngẫu nhiên 1 viên bi. Xét tính đúng sai:
\choiceTF
{\True Khả năng lấy được bi xanh là cao nhất.}
{\True Khả năng lấy được bi vàng là thấp nhất.}
{Chắc chắn sẽ lấy được viên bi màu đỏ.}
{\True Có thể lấy được viên bi màu xanh, đỏ hoặc vàng.}
\loigiai{
\begin{itemize}
    \item 5 xanh (nhiều nhất): Đúng.
    \item 2 vàng (ít nhất): Đúng.
    \item Không chắc chắn lấy được bi đỏ (vì có thể lấy bi khác). Sai.
    \item Có thể ra xanh, đỏ, vàng: Đúng.
\end{itemize}
\begin{itemize}
    \item a) \textbf{Đúng}.
    \item b) \textbf{Đúng}.
    \item c) \textbf{Sai}.
    \item d) \textbf{Đúng}.
\end{itemize}
}
\end{ex}

\textbf{Phần 3. Trả lời ngắn}
\setcounterpageref{ex}{0}

\begin{ex} %[4X9H51]
Tung một đồng xu 25 lần, có 13 lần xuất hiện mặt ngửa. Hỏi mặt sấp xuất hiện bao nhiêu lần?
\shortans{12}
\loigiai{
\textbf{Đáp số:} $12$\\
$25 - 13 = 12$ lần.
}
\end{ex}

\begin{ex} %[4X9H51]
Gieo một con xúc xắc 40 lần, thấy mặt 1 chấm xuất hiện 8 lần. Hỏi các mặt còn lại xuất hiện tất cả bao nhiêu lần?
\shortans{32}
\loigiai{
\textbf{Đáp số:} $32$\\
$40 - 8 = 32$ lần.
}
\end{ex}

\begin{ex} %[4X9V51]
Quay vòng quay 50 lần, mũi tên chỉ vào ô màu đỏ 18 lần, ô màu xanh 22 lần, còn lại là ô màu vàng. Hỏi mũi tên chỉ vào ô màu vàng bao nhiêu lần?
\shortans{10}
\loigiai{
\textbf{Đáp số:} $10$\\
$50 - (18 + 22) = 10$ lần.
}
\end{ex}

\begin{ex} %[4X9V51]
Tung đồng xu 60 lần, số lần xuất hiện mặt ngửa nhiều hơn mặt sấp là 6 lần. Tính số lần xuất hiện mặt ngửa.
\shortans{33}
\loigiai{
\textbf{Đáp số:} $33$\\
$(60 + 6) : 2 = 33$ lần.
}
\end{ex}

\begin{ex} %[4X9V51]
Tung đồng xu 60 lần như trên, tính số lần xuất hiện mặt sấp.
\shortans{27}
\loigiai{
\textbf{Đáp số:} $27$\\
$60 - 33 = 27$ lần.
}
\end{ex}

\begin{ex} %[4X9C51]
Một hộp có 20 thẻ đánh số từ 1 đến 20. Rút ngẫu nhiên 1 thẻ. Có bao nhiêu khả năng rút được thẻ mang số chia hết cho 5?
\shortans{4}
\loigiai{
\textbf{Đáp số:} $4$\\
Các số: 5, 10, 15, 20 (có 4 thẻ).
}
\end{ex}
